\runninghead{Interleaved Commentary}
\section{Exegetical Questions Across the Propers}
\zlabel{triptych:concise:commentary:start}%

\subsection{Why Does God Seem to Harden Hearts and Hide His Face? (Isaiah~63:17, 19b; 64:2b, 4--7 [Vg.~63:17; 64:1, 3b, 5--8]; Psalm~80[79])}

Isaiah opens with a bold complaint: \emph{Quare errare nos fecisti, Domine, de viis tuis; indurasti cor nostrum ne timeremus te?} (``Why hast thou made us to err, O Lord, from thy ways: why hast thou hardened our heart, that we should not fear thee?'', Isa.~63:17). Neither St.~Jerome (\work{Commentarii in Isaiam} XVII, PL~24 cols.~642--651) nor St.~Thomas Aquinas (\work{Expositio super Isaiam ad litteram}, ch.~63--64) makes God the author of sin, yet they diagnose the cause differently.

Jerome, followed by \emph{Forbearance and Unveiling}, traces Israel's wandering and hardness to God's \emph{nimia patientia} (``excessive forbearance''). By holding back immediate punishment to await repentance --- like a physician hiding the scalpel or a father withholding the rod (Jerome cites Hos.~4:14, 2:14; Ps.~89[88]:33; Heb.~12) --- God's slowness looks from below like indifference and lets unchastened hearts grow cold. Aquinas, followed by \emph{Clay, Vine, and Fidelity}, cites Romans~9:18 and explains the verse through permissive causality: God makes Israel err and hardens the heart \emph{permittendo, non immittens gratiam, et non corrigens} (``by permitting, not infusing grace, and not correcting''). St.~Augustine (\work{De gratia et libero arbitrio}, ch.~45 [XXIII]) states the doctrinal boundary beneath both accounts on Exodus~4--14 and Romans~9:18: withholding assisting grace in just judgment never abolishes free will, for ``both God hardened him by His just judgment, and Pharaoh by his own free will.'' \emph{Vigilant Stewardship} takes Isaiah~63:17 with 64:5--7~Vg.\ (Heb.~64:4--6) as a portrait of spiritual sleep --- a people drifting from God's ways and swept away like fallen leaves while no one rises to take hold of him.

In the next selected verses, Jerome and Aquinas diverge sharply again on the sudden past tense of Isaiah~64:3b~Vg.\ (Heb.~64:2b; \emph{descendisti, et a facie tua montes defluxerunt}, ``thou didst come down, and at thy presence the mountains melted away''), standing directly after the wish of 64:1~Vg.\ (Heb.~63:19b; \emph{Utinam dirumperes caelos, et descenderes}, ``O that thou wouldst rend the heavens, and wouldst come down''). Jerome, followed by \emph{Forbearance and Unveiling}, hears the Hebrew tense shift as a prayer granted on the spot, placing second-person address on Israel's lips: \emph{Descendisti, Verbum caro factum es, et habitasti in nobis, vere Emmanuel} (``Thou didst come down; thou, the Word, wast made flesh and didst dwell among us, truly Emmanuel,'' PL~24 col.~646), with the melting ``mountains'' (Deut.~4:24; Ps.~97[96]:5) signifying hostile powers and proud heights raised against God.

Aquinas, followed by \emph{Clay, Vine, and Fidelity}, recognizes that Christological sense --- on 64:1~Vg.\ he explains \emph{dirumperes} (``rend'') as Christ emptying himself, laying aside majesty, and taking flesh (\emph{exinanires temetipsum, deponens majestatem, et carnem assumeres}), and on 64:3b~Vg.\ notes both Final Judgment (\emph{ad judicandum nos}, ``to judge us'') and the Virgin's womb (\emph{in uterum virginis}). Yet for 64:3b~Vg.\ he explicitly prefers ``better, according to the literal sense'' (\emph{vel melius, secundum litteralem sensum}). Having linked omitted 64:3a~Vg.\ (Heb.~64:2a; \emph{non sustinebimus}, ``we shall not bear them'') to Sinai (Exod.~20), Aquinas takes \emph{Descendisti} (``thou didst come down'') as proof from Israel's past: God already came down to deliver Israel from Egypt (Exod.~4; Ps.~18[17]:10) and melted the cliffs of Arnon (Num.~21). Following Jerome, the first interpretation hears the Incarnation already fulfilling Israel's cry; following Aquinas, the second grounds Israel's hope in proven covenant deliverance before its fulfillment in Christ.

Both Doctors identify Isaiah~64:5~Vg.\ (Heb.~64:4; \emph{Ecce tu iratus es, et peccavimus}, ``behold thou art angry, and we have sinned'') as \emph{ordo praeposterus} (inverted order): God was angry because Israel sinned, not the reverse. Aquinas links ``him that rejoiceth, and doth justice'' mystically to Simeon in the Temple (Luke~2), while \emph{Vigilant Stewardship} pairs that clause with the petition of the ancient Gelasian prayer (II.lxxxiv, Postcommunion prayer~683) to go forth fittingly in righteous deeds (\emph{in operibus iustis aptos occurrere}). On verse~6~Vg.\ (Heb.~v.~5; ``all our justices as the rag of a menstruous woman''), Jerome and Aquinas divide a third time. Jerome applies the simile to legal observance after Christ's coming --- unclean beside evangelical purity, as Paul counted legal attainments loss (Phil.~3:8) and Esther loathed her diadem (Esther~14:16). Aquinas instead protects graced virtue by naming four concrete stains on righteous deeds: mortal sin (Sir.~47), negligence (Jer.~48), twisted intention (Matt.~5), and pride (Isa.~57).

At verse~7~Vg.\ (Heb.~v.~6; ``thou hast hid thy face from us,'' \emph{Abscondisti faciem tuam a nobis}), Jerome exhorts sinners to confess and immediately quotes this Sunday's Responsorial Psalm refrain, Psalm~80(79):4 (\emph{Ostende faciem tuam, et salvi erimus}, ``Show thy face, and we shall be saved,'' PL~24 col.~650). In verse~8~Vg.\ (Heb.~v.~7) the speakers abandon self-merit: ``we are clay: and thou art our maker'' (\emph{nos vero lutum; et fictor noster tu}). Jerome notes that looking at human deserts breeds despair whereas appealing to God's Fatherhood (Exod.~4:22) and Potter's craft sustains prayer; Aquinas pairs the verse with Jeremiah~18 and Job~10.

In Psalm~80(79), St.~Augustine (\work{Enarrationes in Psalmos}, Ps.~79 §§1--3, 9, 11) declares that ``The song here is of the Advent of the Lord and of our Saviour Jesus Christ, and of His vineyard... Head and Body, King and people, Shepherd and flock.'' Because \emph{Cherubim} means ``fulness of knowledge'' and love fulfills the Law (Rom.~13:10), any believer holding fast to charity becomes God's throne. On ``Stir up thy might'' (v.~3b), Augustine thinks of the Cross, where Christ seemed weak like the angel who let Jacob prevail so as to bless him. Asking how God's face could ever grow dark (v.~4), Augustine gives the exposition of hiddenness that \emph{Forbearance and Unveiling} follows: Christ placed ``a cloud of flesh, and as it were a veil of weakness'' before his majesty until the Resurrection and Ascension unveiled who he was.

\emph{Clay, Vine, and Fidelity} follows Augustine (Ps.~79 §§9, 11) and St.~Robert Bellarmine (\work{Explanatio in Psalmos}, Ps.~79, explanations~1--3, 14--18) on verses~15--19. Bellarmine maps the omitted tribal names across the promised land (Samaria in Ephraim, Jerusalem in Benjamin, trans-Jordan in Manasseh) and connects ``thou that sittest upon the Cherubim'' to the ark's mercy seat. On verses~15--19 he argues that ``All agree that this passage refers literally to Christ,'' because after the Babylonian captivity the Jewish people never recovered their former position, so the vineyard reached perfection only in Christ --- ``the man of thy right hand'' formed in the Virgin through the Holy Spirit (Luke~1:35). Augustine agrees that God plants no rival vine, perfecting Abraham's seed upon Christ, the sole foundation (1~Cor.~3:11); converted by grace, the people love God \emph{gratis} (for himself alone) and vow ``we will call upon thy name'' (v.~19). Bellarmine also connects ``visit this vineyard'' (v.~15) to Zechariah's canticle (Luke~1:68) and Christ weeping over Jerusalem for missing ``the time of thy visitation'' (Luke~19:44), a warning \emph{Vigilant Stewardship} pairs with the psalm's pledge, ``And we depart not from thee'' (v.~19).

\subsection{How Can Paul Call a Divided Church Enriched and Blameless? (1~Corinthians~1:3--9 and Isaiah~64:3 [Vg.~64:4])}

Paul's gratitude in 1~Corinthians~1:3--9 stands directly ahead of his rebuke of Corinthian factions (1:10--12) and immaturity (3:1--3). Why thank God that a quarrelsome community is ``made rich in him, in all utterance and in all knowledge'' (1:5) and lacks no grace (1:7), and does promising that faithful God (\emph{Fidelis Deus}) will confirm them ``unto the end without crime'' (1:8--9) guarantee perseverance unconditionally?

On verse~7, St.~John Chrysostom (\work{Homilies on First Corinthians}, Homily~II, §6) fastens on one Greek noun: the Corinthians await the \emph{apokalupsis} (Latin \emph{revelatio}, ``unveiling'') of Christ, ``implying that although He be not seen, yet He is, and is present even now, and then shall appear.'' \emph{Forbearance and Unveiling} takes its Advent horizon from that point: believers wait for a veil to lift over a Lord already sustaining his Church.

\emph{Clay, Vine, and Fidelity} follows Chrysostom (Homily~II, §§1--8) and Cornelius a Lapide (\work{Commentaria in omnes divi Pauli epistolas}, Vol.~7, on 1~Cor.~1:3--9) through the tension between 1:4--9 and 1:10. On verses~4--5, the two differ in emphasis. Chrysostom accounts for Paul's praise in three ways: the Corinthians showed early zeal before growing lax; Paul praises the faithful remnant (such as Stephanas's household) without shaming offenders at the start; and Paul uses pastoral management (\emph{oikonomikos}), tying every gift to ``the grace of God'' (\emph{in gratia Dei}) to deflate vanity. Across verses~1--10, Paul ``fasteneth them as with nails to the Name of Christ,'' working ``as it were to arouse those who were heavy-headed after some debauch,'' and applying that Name to ``foment their inflammation, and purge out the corruption of the disease'' (§7). Lapide softens Chrysostom's irony, taking Paul's thanksgiving more directly of genuine spiritual wealth distributed across the community.

On verses~8--9, however, Chrysostom and Lapide converge against fatalism. Chrysostom observes that promising to ``confirm'' them hints quietly that they are wavering; God brings believers to the promised goods ``if we be not quite impatient of His rein,'' quoting Paul's own words from 1~Corinthians~2:9 (``which eye hath not seen, nor ear heard,'' in Paul's eye-before-ear order, recalling Isaiah~64:4~Vg.\ [Heb.~64:3]). Lapide states the same balance through the Council of Trent (Session~VI, ch.~11 and~13, following Augustine): God confirms believers ``so far as His part is'' and will not desert them unless he is first deserted, excluding despair without licensing presumptuous certainty apart from special revelation.

That quotation of 1~Corinthians~2:9 recalls Isaiah~64:4~Vg.\ (Heb.~64:3; ``the eye hath not seen, O God, besides thee, what things thou hast prepared for them that wait for thee''). On Isaiah~64:4~Vg., Jerome defends Paul's paraphrase as a faithful sense-for-sense rendering (\emph{non verbum ex verbo reddens... sed sensuum exprimens veritatem}, ``not rendering word for word, but expressing the truth of the meaning''), rejecting apocrypha --- the \emph{Ascension of Isaiah} (\work{Ascensio Isaiae}) and the \emph{Apocalypse of Elijah} (\work{Apocalypsis Eliae}) --- and Gnostic fabrications (naming Basilides and citing St.~Irenaeus, \emph{Against Heresies} [\work{Adversus Haereses}], Book~I). Aquinas, pairing 1~Corinthians~2:9 with Isaiah~64:4~Vg., lists four unseen gifts: eternal rest (John~14), the kingdom of eternal dignity (Matt.~25), divine refreshment (Ps.~23[22]), and the lamp of eternal light (Ps.~132[131]). Finally, \emph{Vigilant Stewardship} reads Chrysostom's warning about a wavering, factious church beside Mark~13:34: an equipped household cannot afford internal division while awaiting the Day of Christ.

\subsection{Why Withhold the Hour, and Who Must Keep Watch? (Mark~13:33--37 and Omitted 13:32)}

Patristic and medieval commentators read Mark~13:33--37 out of omitted verse~32 (``neither the angels in heaven, nor the Son, but the Father''), confronting three questions: how the Son can speak of an unknown hour, how the householder has left for a far country, and how one doorkeeper (\emph{janitor}) and four night watches bind believers of every century.

On the withheld hour and departure (Mark~13:32--34), \emph{Forbearance and Unveiling} follows St.~Hilary of Poitiers (\work{De Trinitate} IX.58--68, especially §§65 and~67) and St.~Bede (\work{In Marci Evangelium Expositio} IV.13). Answering Arians who inferred creaturely ignorance from Mark~13:32 and Matthew~24:36, Hilary asks how the Lord in whom all wisdom is hidden (Col.~2:3) could fail to know his own coming. Scriptural ``not knowing'' marks an economy (\emph{dispensatio}) of silence. Pairing that silence with Christ's Matthean watchfulness sayings (Matt.~25:13 in IX.65; Matt.~24:44, 46 in IX.67), Hilary argues in IX.67 that Christ ``withholds from us the security of certain knowledge, that our minds may be kept on the stretch by the uncertainty of suspense... for the security of knowledge might breed negligence of the faith.'' St.~Thomas Aquinas later applies Hilary's reasoning directly to Mark~13:33 in the \work{Catena Aurea in Marcum} (13:32--37).

Bede (IV.13) keeps Hilary's line --- noting Jesus says \emph{nescitis} (``you do not know''), not ``we do not know,'' to hold disciples in \emph{pendulae exspectationis incerto} (``the uncertainty of poised expectation'') --- and explains the householder's journey by drawing on St.~Gregory the Great's Homily~9 on Matthew~25:14--30 (\work{Homiliae in Evangelia}, Hom.~9): because earth is the native home of human flesh, Christ traveled into a ``far country'' when he carried risen humanity to heaven at the Ascension. He left his Church bodily (\emph{corporaliter}), without withdrawing his divine safeguard.

On the next phrase --- \emph{dedit servis suis potestatem cujusque operis} (``gave authority to his servants over every work,'' Mark~13:34) --- \emph{Clay, Vine, and Fidelity} adopts Bede's identification of that authority with the grace of the Holy Spirit empowering every good work. Mark's charge to watch and pray rests on prior endowment rather than unaided willpower, though Mark~13's abrupt eschatological warning is also the main difficulty for the second interpretation, whose Pauline and Isaian grace theology can mute the Gospel's alarm if held apart from vigilance.

\emph{Vigilant Stewardship} puts Mark~13:34--37 in the foreground, keeping four distinct patristic arguments clear:
\begin{itemize}[leftmargin=1.5em,itemsep=0.2em]
\item \textbf{Pastors and laity:} Bede (IV.13) identifies the \emph{janitor} with pastors and spiritual rulers guarding the flock, while verse~37 (\emph{omnibus dico}, ``I say to all'') requires every Christian to guard the door of his own heart against evil suggestions. Quoting Gregory's Homily~13 (on Luke~12:35--40), Bede defines watching concretely: \emph{Vigilat qui ad aspectum veri luminis mentis oculos apertos tenet, vigilat qui servat operando quod credit, vigilat qui a se torporis et negligentiae tenebras repellit} (keeping the mind's eyes open to the true light, keeping in deeds what one believes, and driving off torpor and negligence).
\item \textbf{Watches as ages of life:} Gregory's Homily~13 maps the three Lucan night watches (Luke~12:38) onto boyhood, youth, and old age, portraying God as a \emph{patiens redditor} (``patient judge'') who leaves no age hopeless yet demands stricter account if his patience is spurned, keeping the hour hidden \emph{ut semper possit esse suspecta} (``so that it may always be anticipated''). Mediated through Aquinas's \work{Catena Aurea in Marcum} (13:32--37), Theophylact of Ohrid distributes Mark~13:35's four Roman watches differently (evening = old age, midnight = mid-youth, cockcrow = mature reason, morning = childhood) and warns that wicked men also stay awake at night, which is why Christ joins watching with prayer.
\item \textbf{Sleep and a fading world:} In the same \work{Catena Aurea} passage, Pseudo-Jerome likens spiritual sleep to chasing dream-phantoms that vanish upon waking --- echoing Gregory's Homily~1 (on Luke~21:25--33, §§1 and~5), which depicts the world growing old under affliction so believers stop clinging to what cannot last.
\item \textbf{Every century bound:} Through the \work{Catena Aurea in Marcum}, Aquinas preserves St.~Augustine's solution in \work{Epistula}~199 §3 to the objection that only one generation will see the Last Day: ``that day comes to each man when his day comes for departing from this life such as he is to be, when judged in that day.'' Because death fixes the state in which each believer meets the Judge, every generation stands equally under the command to watch.
\end{itemize}

\subsection{How Do the Two Advents and the Altar Shape Prayer and Daily Life? (Antiphons, Orations, Preface, and Communion)}

The Missal's shared Advent texts receive distinct emphases across the three interpretations. In Psalm~25(24), beyond verses~1--3 at the Entrance, Augustine (\work{Enarr.}, Ps.~24 §10) identifies ``all the ways of the Lord are mercy and truth'' with Christ's two Advents --- ``the one in mercy, the other in judgment,'' forgiving sins without desert at the first and testing works in truth at the second. Reading the ancient Gelasian Postcommunion prayer (\work{Gelasianum Vetus} II.lxxxiv, no.~683, whose exact lexical alignment with the modern Collect was not collated offline), \emph{Clay, Vine, and Fidelity} stresses its petition asking God for the very \emph{voluntatem} (``will'') to meet Christ, \emph{Vigilant Stewardship} stresses going forth \emph{in operibus iustis} (``in righteous deeds,'' recalling Isa.~64:5~Vg.\ [Heb.~64:4]), and \emph{Forbearance and Unveiling} takes that petition as active readiness during delay.

At the Gospel Acclamation (\emph{Ostende nobis, Domine}, Ps.~85[84]:8), Augustine (\work{Enarr.}, Ps.~84 §6) identifies ``thy salvation'' as ``Thy Christ'' and turns the petition inward against pride: knowing that every praised good comes from mercy keeps the believer humble, purifying the heart to see Christ ``not as the Jews saw Him and crucified Him, but as the Angels see Him, and rejoice.'' Bellarmine (\work{Explanatio}, Ps.~84, explanation~7) anchors the plea in expecting the Incarnation (Titus~2:11; 3:4) while endorsing Augustine's moral counsel.

At the altar, the ancient \work{Veronense} prayer (XXVIIII.372, historical counterpart to the modern Prayer over the Offerings, \emph{Suscipe, quaesumus, Domine}) asks both for grace and for faithful offering: \emph{Clay, Vine, and Fidelity} emphasizes that the Church brings back gifts drawn from God's own benefits (\emph{de tuis... collata beneficiis}), while \emph{Vigilant Stewardship} and \emph{Forbearance and Unveiling} emphasize consecrating temporal devotion (\emph{devotioni... temporalis}) toward eternal redemption. At the Seasonal Preface (GIRM~364; the modern Preface wording is not printed here), three teachers expound Christ's comings in ways the three interpretations draw on:
\begin{itemize}[leftmargin=1.5em,itemsep=0.2em]
\item \textbf{St.~Cyril of Jerusalem (\work{Catechesis}~15, §§1--2, 4, 26):} ``We preach not one advent only of Christ, but a second also, far more glorious than the former.'' Christ has a twofold generation and twofold descents --- ``one, the unobserved, like rain on a fleece; and a second His open coming, which is to be,'' moving from swaddling clothes, manger, and Cross to light, angels, and the crown of a divine kingdom (§1; Mal.~3:1--2 in §2). Cyril forbids date-curiosity (Acts~1:7, §4) and sums up readiness from Matthew~25:35--41 (§26): ``there is here no need of allegory, but to fulfil what is said'' in works of mercy.
\item \textbf{St.~Bernard of Clairvaux (\work{Sermones in Adventu}~3.4, 5.1):} Between the first coming ``in the flesh and in weakness'' and the last ``in glory and majesty,'' Bernard teaches a middle coming \emph{in spiritu et virtute} (``in spirit and in virtue''). Distinguishing three comings --- \emph{ad homines, in homines, contra homines} (``to men, into men, against men'') --- he cites John~14:23 in Sermo~3.4 on Christ dwelling invisibly in those who love him and keep his word, and in Sermo~5.1 expounds the psalmist's promise in Psalm~68(67):14 to ``sleep among the midst of lots,'' an image of hidden nearness that \emph{Forbearance and Unveiling} follows.
\item \textbf{St.~Leo the Great (\work{Sermon}~19, §§1--3):} Preaching the December fast from Luke~21:34 because ``the day denounced is near, even though hidden'' (§1), Leo sets out the Roman Advent discipline that \emph{Vigilant Stewardship} adopts: pairing bodily fasting with a fast of the mind from vice (§2) and generous almsgiving to the poor (§3).
\end{itemize}

At Holy Communion (\emph{Dominus dabit benignitatem}, Ps.~85[84]:13), Bellarmine (\work{Explanatio}, Ps.~84, explanation~12) and Augustine (\work{Enarr.}, Ps.~84 §§10--12) offer complementary expositions. Bellarmine identifies ``our earth'' with the untouched Virgin Mary overshadowed by the Spirit to bring forth Christ like the ``flower of the field'' (Song~2:1). Augustine expounds verse~12 (``Truth hath sprung out of the earth'') both of Christ taking flesh from Mary for the Cross (§10) and of sinners confessing truth (§11); on verse~13 (§12), rendering \emph{benignitas} as ``sweetness,'' he explains that God gives the interior delight of working righteousness, transforming tastes so hearts ``furrowed, as it were, with the plough of confession'' delight in sobriety, giving, prayer, and hymns rather than drunkenness, theft, or idle shows. At the altar, Christ the fruit of Mary enters communicants in hidden nearness (\emph{Forbearance and Unveiling}), enables human clay to bear fruit (\emph{Clay, Vine, and Fidelity}), and sustains concrete temperance and almsgiving (\emph{Vigilant Stewardship}); the Prayer after Communion (\emph{Prosint nobis}) closes the Communion Rite, and the Concluding Blessing precedes the dismissal (GIRM~90).

\subsection{Comparing the Three Interpretations and Their Four Senses}

Genuine exegetical disagreements separate the Fathers and Doctors on three points: Jerome's \emph{nimia patientia} and Incarnational \emph{Descendisti} (Isa.~63:17; 64:3b~Vg.\ [Heb.~64:2b]) against Aquinas's permissive non-infusion of grace and literal Exodus/Arnon deliverance; Jerome's anti-legalist exposition of Isaiah~64:6~Vg.\ (Heb.~64:5) against Aquinas's four concrete stains on righteous acts; and Chrysostom's ironic pastoral account of 1~Corinthians~1:4--5 against Lapide's more direct account of Corinthian gifts (alongside Gregory's three Lucan watches and Theophylact's four Markan watches). Other pairings --- Augustine's moral sweetness and Bellarmine's Marian fruit on Psalm~85(84):13, or Hilary's economy of suspense, Bede's clerical and lay doorkeeping, and Augustine's personal-death horizon on Mark~13 --- complement one another.

Those choices yield three distinct four-sense profiles that check one another's blind spots. Under \emph{Forbearance and Unveiling}, the literal sense begins with perceived divine silence, the allegorical dwells on the veiled Incarnation and bodily Ascension, the moral urges patient prayer against numbness, and the anagogical awaits the \emph{apokalupsis} of an already-present Lord; pressed alone, it risks softening the guilt confessed in Isaiah~64:5--7~Vg.\ (Heb.~64:4--6) and 1~Corinthians~1:10's warning against factions. Under \emph{Clay, Vine, and Fidelity}, the literal sense rests on God's Exodus deliverance (Aquinas on Isa.~64:3b~Vg.\ [Heb.~64:2b]) and covenant vine, the allegorical rests on Christ the Son of man born of Mary and giving the Spirit's authority (Bede on Mark~13:34), the moral calls for confession and willing cooperation with grace, and the anagogical trusts \emph{Fidelis Deus} to confirm persevering believers blameless; its hazard is letting grace theology mute Mark~13's urgent alarm. Under \emph{Vigilant Stewardship}, the literal sense centers on Jesus' parting charge on the Mount of Olives, the allegorical presents pastors and faithful living between the two Advents, the moral presses daily fidelity, temperance, and almsgiving across every age, and the anagogical sets personal death beside the universal Parousia; detached from Isaiah~64:6--8~Vg.\ (Heb.~64:5--7) and 1~Corinthians~1:4--9, active service can drift into self-reliant moralism or leave bewildered sufferers unconsoled.

\begin{comparisontable}{Dimension}{1. Forbearance and Unveiling}{2. Clay, Vine, and Fidelity}{3. Stewardship Between Advents}
Starting question & Why does God seem absent or silent while hearts harden and the master tarries? & How can stained, divided believers hope to stand blameless before Christ? & Who is charged to watch in Mark~13, and what concrete life counts as watching? \\
Key textual cruxes & Isa.~63:17 (\emph{nimia patientia}), 64:1, 3b~Vg.\ (\emph{Descendisti}), 1~Cor.~1:7 (\emph{apokalupsis}), Mark~13:32--34 (withheld hour; bodily departure) & Isa.~63:17 (permission), 64:3b~Vg.\ (Egypt and Arnon), 64:6--8~Vg.\ (four stains; Potter), 1~Cor.~1:4--9 (\emph{Fidelis Deus}), Ps.~80(79):15--19 & Mark~13:34--37 (doorkeeper vs.\ all; four watches; personal death), Ps.~25(24):10 and Cyril \work{Cat.}~15 (two Advents), Isa.~64:5~Vg. \\
Principal Fathers and Doctors & Jerome (Isa.~63--64), Chrysostom (1~Cor.~1:7), Hilary and Bede (Mark~13:32--34) & Aquinas (Isa.~63--64), Augustine (hardening; Pss.~25[24], 80[79], 85[84]), Bellarmine (Pss.~80[79], 85[84]), Chrysostom (1~Cor.~1:4--9) & Bede (Mark~13:34--37), Gregory (Hom.~1, 9, 13), Augustine (Ps.~25[24]:10; \work{Ep.}~199), Cyril (\work{Cat.}~15), Leo (\work{Sermon}~19) \\
Isa.~63:17 and 64:3b~Vg.\ (Heb.~64:2b) & Jerome: forbearance; Incarnation directly in \emph{Descendisti} & Aquinas: permissive non-infusion; Egypt (Exod.~4) and Arnon (Num.~21) literally & Portrait of spiritual drowsiness set against ``him that... doth justice'' (\emph{facienti justitiam}, Isa.~64:5~Vg.) \\
Pastoral emphasis & Consoling and contemplative; addresses bewilderment at divine silence & Humbling and reassuring; anchors perseverance in grace and willing cooperation & Mobilizing and vocational; guides clergy and laity across every age of life \\
Risk to guard against & Under-weighting human guilt and sin in Isa.~64:5--7~Vg.\ (Heb.~64:4--6) & Letting grace theology mute Mark~13's urgent watchfulness & Sliding into self-reliant activism detached from prior grace \\
\end{comparisontable}
